\documentclass[10pt,a4paper]{amsart}
\usepackage[margin=19mm]{geometry}
\usepackage{amsmath,amssymb,mathtools}
\usepackage[scr=rsfs,cal=euler]{mathalfa}
\usepackage{xfrac}
\usepackage[unicode,hidelinks,bookmarks=false]{hyperref}
\newcommand{\ie}{i.e.\ }
\newcommand{\cf}{cf.\ }
\newcommand{\resp}{respectively\ }
\newcommand{\Subsection}{\subsection}
\providecommand{\nobreakdash}{\nobreak\hbox{-}}
\setlength{\emergencystretch}{2em}
\multlinegap0pt
\usepackage{bm}
\usepackage{bbm}
\usepackage{dsfont}
\usepackage{xspace}
\usepackage{cancel}
\usepackage{mathtools}
\usepackage{amsthm}
\makeatletter
\@ifundefined{thm}{\newtheorem{thm}{Theorem}}{}
\@ifundefined{lem}{\newtheorem{lem}[thm]{Lemma}}{}
\makeatother
\title[Sign games on graphs]{Sign games on graphs}
\author{Liz Blum, Lily Brustkern, Rosetta Hawkins, Neil R. Nicholson, Ranjan Rohatgi}
\date{Source: arXiv:2511.05451v2 (CC BY 4.0)}
\begin{document}
\maketitle

\noindent\emph{Source and licence.} Sign games on graphs. Version arXiv:2511.05451v2 (CC BY 4.0), distributed under the Creative Commons Attribution 4.0 International licence, as recorded in the arXiv licence metadata for this exact version and shown on the abstract page https://arxiv.org/abs/2511.05451v2. Full rights notice: licenses/arxiv-2511.05451-CC-BY-4.0.txt.\par\smallskip

\noindent\textit{Editorial scope.} Counted unit: the Thm whose source label is \texttt{thm1} in the article (source0.tex lines 198--200, source label \texttt{thm1}), delivered together with the article's own complete original proof of it (source0.tex lines 201--224, one \texttt{proof} environment of 24 lines). No mathematics is written, repaired or re-derived: statement and proof are the article's own text line for line. Required context retained uncounted: lem1 (lines 184--190). Every definition, display and label of those lines is retained unchanged. Omitted from this package and not redistributed: 1--183, 191--197, 225--736. Nothing inside a retained excerpt is omitted or altered: every retained body is delivered exactly as the article's lines are written, so no omission receipt of any kind accompanies this package. Citations: the retained excerpts carry no citation command at all, so no bibliography entry of the article is transcribed and no reference is redistributed. Reference adaptation: none was needed and none was made; every reference inside the delivered unit resolves inside it, so no unresolved reference or citation is produced. Printed slips kept literal: no printed word, symbol, hypothesis or number of the counted statement or of its proof was altered. Layout and class: the article's own class is \texttt{article} with packages including setspace, amsmath, amsthm, amssymb, graphicx, subfigure, epsfig, tikz, hyperref, chngcntr; the wrapper is the lane's pinned \texttt{amsart} editorial wrapper at 10pt on A4, which restates the theorem-like environments the retained text needs and adds the article's own macro definitions that the retained text needs (none), with extra wrapper packages (bm, bbm, dsfont, xspace, cancel, mathtools). The delivered numbering is the unit's own. Assets: no retained span contains a figure environment, a graphics-inclusion command, a PDF-page inclusion, a table or a TikZ picture; no image file is copied into this package, the package declares no asset dependency and no image file is redistributed with it. Licence: version arXiv:2511.05451v2 of this article is distributed under the Creative Commons Attribution 4.0 International licence, as recorded in the arXiv licence metadata for this exact version and shown on the abstract page https://arxiv.org/abs/2511.05451v2. A full rights notice is delivered at licenses/arxiv-2511.05451-CC-BY-4.0.txt. Characters: every retained excerpt is pure ASCII, so the delivered engine, XeLaTeX with its default Latin Modern text font, reports no missing character.
\begin{lem} \label{lem1}
If the Sign Game is played to completion on $G = K_n$ and $a$ vertices are assigned $+1$ and $b$ vertices assigned $-1$ (so that $a+b = n$, with $a,b \geq 2$), then 

\begin{equation}
         s(G) = \binom{a}{2}+ \binom{b}{2}-ab.
    \end{equation}   
\end{lem}

\begin{thm} \label{thm1}
    Player N will win the Sign Game played on $K_n$ for $n\geq 2$, unless Player N is Player 1 and either $n=2$, in which case Player P wins, or $n=4$, which results in a draw.
\end{thm}

\begin{proof}
    As in Lem. \ref{lem1}, suppose the Sign Game is played to completion on $G=K_n$ ($n \geq 2$) with $a$ vertices assigned $+1$ and $b$ vertices assigned $-1$. Without loss of generality, we can assume $a \geq b$, so that $a = b+r$ for some $r \geq 0$. By Lem. \ref{lem1}, we have

    \begin{align}
        s(G) &= \binom{a}{2} + \binom{b}{2} -ab \\
        &= \binom{b+r}{2} + \binom{b}{2}-(b+r)b \\
        &=\frac{r^2-r}{2}-b,
    \end{align}

    
    \noindent which is negative precisely when 

    \begin{equation} \label{completeinequality}
         \frac{r^2-r}{2}<b.
    \end{equation}

    Player N can guarantee $r \leq 2$ by mirroring any turn by Player P with an opposite vertex assignment. Should Player N have the first turn of the game, then any vertex assignment, followed by this strategy, will ensure that $r \leq 2$.

    Note that when $r=0$ or $r=1$, Eqn. \ref{completeinequality} is satisfied so long as $b \geq 1$. The described mirroring strategy yields such a value of $b$ (and consequently $r$) in the cases when Player P has the first turn, or, when $n$ is odd.

    For all other scenarios, if $r=2$, then $s(G) < 0$ when $1 < b$ (that is, Player N will win as long as there are at least two vertices assigned a $-1$). Following the described mirroring strategy, this will occur for all cases except $n =2$ or $n=4$ along with Player N having the first turn.

    So, we must consider these cases: playing on $K_2$ or $K_4$ with Player N being Player 1. The first, on $K_2$, is obvious as Player P can guarantee $s(G)=1$. When $G = K_4$, there are only two possible outcomes (as players are playing optimally) to consider. The first, $a = 3$ and $b =1$ (or equivalently $a = 1$ and $b = 3$) yields $s(G) = 0$. The other possibility of $a = b = 2$ results in a score of $-2$. Player P can guarantee this does not happen by mirroring Player N and assigning vertices identically, meaning either $a$ or $b$ is at least $3$. Thus, playing on $K_4$ results in a draw.
\end{proof}

\end{document}
